\subsubsection*{3.3. 条件 (B)}

\paragraph{3.3.1.}
称包含映射
\begin{equation*}
m:X_{\bar\eta}\hookrightarrow X\times_k k(\bar\eta)
\end{equation*}
满足条件 (B), 若存在
\begin{equation}\tag{3.3.2}
V_{\bar\eta},\quad
\text{为下列空间上余维为 $n-1$ 的代数圈: }
X_{\bar\eta}\times\bigl(X\times_k k(\bar\eta)\bigr),
\end{equation}
使其定义的映射
\begin{equation}\tag{3.3.3}
V_{\bar\eta}:H^{n-1}(X_{\bar\eta},\mathbb Q_\ell(i))
\longrightarrow H^{n-1}(X,\mathbb Q_\ell(i))
\end{equation}
为
\begin{equation}\tag{3.3.4}
m^*:H^{n-1}(X,\mathbb Q_\ell(i))
\longrightarrow H^{n-1}(X_{\bar\eta},\mathbb Q_\ell(i)),
\end{equation}
的左逆, 即满足
\begin{equation}\tag{3.3.5}
V_{\bar\eta}\,m^*=\mathrm{id}.
\end{equation}

\subsubsection*{3.4.}
若 [4, 1.4.2.1] 中的运算 $\wedge$ 是代数的 (参见~[4, 2.12(ii)]), 则条件 (B) 成立. 猜想 $\wedge$ 总是代数的, 因而条件 (B) 总成立; 这在许多情形下已知 [4]. 当 $X$ 是偶数维完全交时, 它平凡成立, 因为此时已知 (IX~1.6)
\begin{equation}\tag{3.4.1}
H^{n-1}(X,\mathbb Q_\ell(i))=0.
\end{equation}
